% !TeX root = ../教材题库_全部条目_转写工作稿.tex
% Source: XJ-B1-C05-HW-P201-Y1548
% 第5章 / 5.5 三角函数模型的简单应用 / 教材 p.201、202 / PDF 第208、209页
% OCR未校对：公式与文字均可能有误；未完成数学TeX转写；禁止直接发给学生
% 原题及相邻原解截图（核对用）：../题目原图/XJ-B1-C05-HW-P201-Y1548-01.webp, ../题目原图/XJ-B1-C05-HW-P201-Y1548-02.webp
% 本文件只在转写完成后作为题干使用；OCR 条目尚不是可靠题干。
\begin{bankitem}{习题5.5 第2题}{第5章 / 5.5 三角函数模型的简单应用 / 教材 p.201、202 / PDF 第208、209页}{XJ-B1-C05-HW-P201-Y1548}{OCR未校对：公式与文字均可能有误；未完成数学TeX转写；禁止直接发给学生}
\begin{ocrdraft}
[教材p.201 / PDF第208页]
\par
如图为某简谐振动的图象，它符合十@(A>O,
\par
2．
\par
2
\par
[教材p.202 / PDF第209页]
\par
<的形式．
\par
（1）求该简谐振动的振幅、周期、频率和初相；
\par
（2）求该简谐振动的函数解析式；
\par
（3）求该函数的单调递增区间．
\end{ocrdraft}
\end{bankitem}
